\documentclass[a4paper,12pt]{article}

\usepackage[a4paper,margin=2cm]{geometry}
\usepackage{amsmath}
\usepackage{amssymb}
\usepackage{graphicx}
\usepackage{enumitem}
\usepackage{titlesec}

% Prevent vertical stretching
\raggedbottom

% Compact spacing
\setlength{\parindent}{0pt}
\setlength{\parskip}{4pt}
\titlespacing*{\section}{0pt}{6pt}{6pt}

\begin{document}

% ================= HEADER =================
\noindent

\includegraphics[width=4cm]{logo.jpeg}

\vspace{0.3cm}

{\large \textbf{IIITB COMET FOUNDATION}} \\ 
\textbf{Name:} Jera Prakash \\
\textbf{Batch:} COMETFWC038 \\
\textbf{Date:} \today

\vspace{0.3cm}
\hrule
\vspace{0.4cm}

% ================= QUESTION =================

\section*{GATE CS 2019 – Question 6 Analysis}

\textbf{Question:}

In 16-bit 2’s complement representation, the decimal number $-28$ is:

(A) 1111 1111 0001 1100  
(B) 0000 0000 1110 0100  
(C) 1111 1111 1110 0100  
(D) 1000 0000 1110 0100  

\vspace{0.2cm}
\hrule
\vspace{0.3cm}

\textbf{Step 1: Convert +28 to Binary}

\[
28_{10} = 11100_2
\]

16-bit form:

\[
0000\ 0000\ 0001\ 1100
\]

\textbf{Step 2: 1’s Complement}

\[
1111\ 1111\ 1110\ 0011
\]

\textbf{Step 3: Add 1}

\[
1111\ 1111\ 1110\ 0011
+
0000\ 0000\ 0000\ 0001
=
\mathbf{1111\ 1111\ 1110\ 0100}
\]

\textbf{Correct Option: (C)}

% ================= HARDWARE =================

\section*{Hardware Implementation}

Verified using Arduino UNO.

\textbf{PlatformIO Code:}

\begin{verbatim}
#include <Arduino.h>
#include <stdint.h>

int16_t num = 28;
int16_t twos_comp;

void setup() {
    twos_comp = -num;
}

void loop() {
}
\end{verbatim}

\textbf{Components Used:}

\begin{itemize}[noitemsep]
\item Arduino UNO
\item Breadboard
\item Common Anode 7-Segment Display
\item 220$\Omega$ Resistor
\item Jumper Wires
\end{itemize}

\textbf{Pin Connections:}

\begin{center}
\begin{tabular}{|l|l|}
\hline
Segment a & Digital 2 \\ \hline
Segment b & Digital 3 \\ \hline
Segment c & Digital 4 \\ \hline
Segment d & Digital 5 \\ \hline
Segment e & Digital 6 \\ \hline
Segment f & Digital 7 \\ \hline
Segment g & Digital 8 \\ \hline
GND & GND \\ \hline
VCC & 5V \\ \hline
\end{tabular}
\end{center}

\textbf{Observation:}

The display showed \textbf{1}, confirming correctness.

% ================= OUTPUT =================

\section*{Hardware Output}

\begin{center}
\includegraphics[height=7.5cm]{output.jpeg}
\end{center}

\vspace{-0.2cm}

\textbf{Conclusion:}

The 16-bit 2’s complement of -28 is:

\[
1111\ 1111\ 1110\ 0100
\]

Hardware and theoretical results match.

\enlargethispage{1cm}

\end{document}
